\documentclass[12pt]{article}

\usepackage[a4paper, total={6in, 8in}]{geometry}
\usepackage{textcomp}

\begin{document}
\setlength{\parindent}{0pt}

\def \t {\textrightarrow}
\def \v {\vspace{1em}}
\def \bi {\begin{itemize}}
\def \ei {\end{itemize}}
\def \s[#1] {\section*{#1}}
\def \ss[#1] {\subsection*{#1}}
\def \sss[#1] {\subsubsection*{#1}}

Even young men from the colonies tried to be involved

\ss[The empire needs men]
Propaganda for men from the colonies \t symbolism of the young lions (colonies) and old lions (uk) \\
In this propaganda symbolism is used a lot \\

\ss[Women of britain say work]
\ss[Isnt all work]
Men would have had recreational activities too \t not only fighting

\v

Soldiers in WWI did not spend most of their time fighting, but in trenches = terrible living conditions \\
When they fought, they were fighting against distruptive new war machines (not other soldiers) \\
They left with some expectations, but they found a completely different situation \\
Many suffered from PTSD \t these soldiers are called shell-shock \\
"To go over the top" = to attack the nearest trench \t "no man's land" is the area separating two trenches

\ss[Maggio 13]
PTSD = post traumatic stess disorder \t this form was called "shell-shocked", and they are called "shell-shocked soldiers" \t and they were not traeted \\
In case of severe shell-shock they were recovered in the "craiglockhart hospital" \t they were sent here not in order to let them adjust to civilian life again, but to treat them to send them again to the front \\
Some survived and return to the UK as war veterans, who after the end of the great war had a hard time in adjusting to civilian life again \\
British women also played a central role in the british war efforts \t they were not soldiers, but they contributed:
\bi
  \item by replacing men in factories
  \item by joining the propaganda system
  \item they were sent to the front as nurses \t there is a lot of british literature of female authors that were nurses in WWI
\ei
After the great war women were rewared by being granted the right to vote (in 1919) \t initially only women over 30 could vote, now over 18 \\
War poets = professional poets who ended up in war
\end{document}
